\oldpage[IV-2]{143}

\subsection{\texorpdfstring{Planitud y propiedad $(R_k)$.}{Planitud y propiedad (R\string_k)}}

\begin{proposition}[6.5.1]
\EGAReaderTarget{ega4.c04.prop.6-5-1}{IV.2.s6.5.1-fr}
Sean $A$, $B$ dos anillos locales noetherianos, $k$ el cuerpo residual
de $A$ y $\varphi:A\to B$ un homomorfismo local para el cual $B$ es un
$A$-módulo plano. Entonces:
\begin{enumerate}
\item[\textup{(i)}] Si $B$ es regular, $A$ es regular.
\item[\textup{(ii)}] Si $A$ y $B\otimes_Ak$ son regulares, $B$ es regular.
\end{enumerate}
\end{proposition}

Esta proposición se da a título de recordatorio, pues ya ha sido demostrada en
(\textbf{0}, 17.3.3).

\begin{corollary}[6.5.2]
\label{IV.2.s6.5.2-fr}
Sean $X$, $Y$ dos preesquemas localmente noetherianos $f:X\to Y$
un morfismo plano.
\begin{enumerate}
\item[\textup{(i)}] Si $X$ es regular en un punto $x$, $Y$ es regular en el punto
$f(x)$.
\item[\textup{(ii)}] Si, para un $y\in f(X)$, $Y$ es regular en el punto $y$ y si
$f^{-1}(y)$ es un preesquema regular en un punto $x$, entonces $X$ es
regular en el punto $x$.
\end{enumerate}
\end{corollary}

\begin{proposition}[6.5.3]
\label{IV.2.s6.5.3-fr}
Sean $X$, $Y$ dos preesquemas localmente noetherianos y $f:X\to Y$
un morfismo plano.
\begin{enumerate}
\item[\textup{(i)}] Si $X$ posee la propiedad $(R_k)$ en un punto $x$
\textup{(resp.} si $X$ posee la propiedad $(R_k)$), $Y$ posee la propiedad
$(R_k)$ en el punto $f(x)$ \textup{(resp.} $Y$ posee la propiedad $(R_k)$ en todos
los puntos de $f(X)$).
\item[\textup{(ii)}] Supongamos que, para todo $y\in f(X)$, el preesquema
$f^{-1}(y)$ posee la propiedad $(R_k)$; entonces, si para un punto
$y\in f(X)$, $Y$ posee la propiedad $(R_k)$ en el punto $y$, $X$ posee
la propiedad $(R_k)$ en todo punto de $f^{-1}(y)$.
\end{enumerate}
\end{proposition}

\begin{proof}
\medskip\noindent
\begin{enumerate}
\item[(i)] Pongamos $y=f(x)$ y sea $y'$ una generalización de $y$; como
en la demostración de (6.4.1), existe un punto genérico $x'$ de una
componente irreducible de $f^{-1}(y')$ que es una generalización de $x$;
por tanto, se tiene
$\dim(\mathcal O_{x'}\otimes_{\mathcal O_{y'}}\boldsymbol{k}(y'))=0$
y en virtud de la hipótesis y de (6.1.2),
$\dim(\mathcal O_{x'})=\dim(\mathcal O_{y'})$.
La hipótesis implica, bien que
$\dim(\mathcal O_{\boldsymbol{x}'})>k$, en cuyo caso
$\dim(\mathcal O_{y'})>k$, bien que $\mathcal O_{x'}$ es un anillo
regular, y entonces $\mathcal O_{y'}$ es un anillo regular en virtud de
(6.5.1).

\item[(ii)] Como, para toda generalización $x'$ de $x\in f^{-1}(y)$,
$y'=f(x')$ es una generalización de $y$, es posible limitarse a verificar que,
si $\dim(\mathcal O_x)\leqslant k$, entonces $\mathcal O_x$ es un anillo
regular. Ahora bien, se tiene en virtud de (6.1.2)
\[
\dim(\mathcal O_x)=\dim(\mathcal O_y)
+\dim(\mathcal O_x\otimes_{\mathcal O_y}\boldsymbol{k}(y))
\]
por tanto, si $\dim(\mathcal O_x)\leqslant k$, se tiene \emph{a fortiori}
$\dim(\mathcal O_y)\leqslant k$ y
$\dim(\mathcal O_x\otimes_{\mathcal O_y}\boldsymbol{k}(y))\leqslant k$ y la hipótesis
implica que $\mathcal O_y$ y
$\mathcal O_x\otimes_{\mathcal O_y}\boldsymbol{k}(y)$ son anillos regulares. Se concluye
entonces de (6.5.1) que $\mathcal O_x$ es un anillo regular.
\end{enumerate}
\end{proof}
